\documentclass[titlepage,12pt,a4paper,times]{book}

\usepackage[utf8]{inputenc}
\usepackage[portuguese]{babel}
\usepackage[T1]{fontenc}
% substituir linha anteior por 
% \usepackage[english]{babel} 
% se o relatório for escrito na língua inglesa.
\usepackage{makeidx}
\usepackage{xspace}
\usepackage{graphicx,color,times}
\usepackage{fancyhdr}
% \usepackage{pxfonts}
% \usepackage{times}
% \usepackage{mathptm}
% \usepackage{amssymb}
% \usepackage{amsfonts}
\usepackage{amsmath}
\usepackage{latexsym}
\usepackage[printonlyused]{acronym}
\usepackage{float}
\usepackage{listings}
\usepackage{tocbibind}
\usepackage{wrapfig}
\usepackage{natbib}
\usepackage{hyperref}
% \usepackage{glossaries}
% \makeglossaries

\renewcommand{\ttdefault}{phv}

\pagestyle{fancy}
\renewcommand{\chaptermark}[1]{\markboth{#1}{}}
\renewcommand{\sectionmark}[1]{\markright{\thesection\ #1}}
\fancyhf{} \fancyhead[LE,RO]{\bfseries\thepage}
\fancyhead[LO]{\bfseries\rightmark}
\fancyhead[RE]{\bfseries\leftmark}
\renewcommand{\headrulewidth}{0.5pt}
\renewcommand{\footrulewidth}{0pt}
\addtolength{\headheight}{0.5pt}
\setlength{\marginparsep}{0cm}
\setlength{\marginparwidth}{0cm}
\setlength{\marginparpush}{0cm}
\addtolength{\hoffset}{-1.0cm}
\addtolength{\oddsidemargin}{\evensidemargin}
\addtolength{\oddsidemargin}{0.5cm}
\addtolength{\evensidemargin}{-0.5cm}


% NEW COLORS
\definecolor{dark}{gray}{0.25}
\definecolor{lgray}{gray}{0.9}
\definecolor{dkblue}{rgb}{0,0.13,0.4}
\definecolor{dkgreen}{rgb}{0,0.6,0}
\definecolor{gray}{rgb}{0.5,0.5,0.5}
\definecolor{mauve}{rgb}{0.58,0,0.82}

\lstset{ %
  language=C,                    basicstyle=\footnotesize,
  numbers=none,                  numberstyle=\tiny\color{gray}, 
  stepnumber=1,                  numbersep=5pt,
  backgroundcolor=\color{white}, showspaces=false,
  showstringspaces=false,        showtabs=false,
  frame=single,                  rulecolor=\color{black},
  tabsize=2,                     captionpos=b,
  breaklines=true,               breakatwhitespace=false,
  title=\lstname,                keywordstyle=\color{blue},
  commentstyle=\color{dkgreen},  stringstyle=\color{mauve},
  escapeinside={\%*}{*)},        morekeywords={*},
  belowskip=0cm
}



\begin{document}


\thispagestyle{empty}
\setcounter{page}{-1}

\begin{center}
\begin{Huge}
\textbf{Universidade da Beira Interior}
\end{Huge}
\end{center}

\begin{center}
\begin{Huge}
Departamento de Informática
\end{Huge}
\end{center}

\vspace{0,07cm}
\begin{figure}[!htb]
\centering
\includegraphics[scale=0.3]{brasaoubi.JPG}
\end{figure}

\vspace{0.5cm}
\begin{center}
\begin{Large}
\textbf{Nº 01 - 2016: \emph{Secure Gateway for Domestic Networks}}
\end{Large}
\end{center}


\vspace{0.5cm}
\begin{center}
\begin{normalsize}
\begin{large}
Authored by:
\end{large}
\end{normalsize}
\end{center}

\vspace{0.2cm}
\begin{center}
\begin{large}
\textbf{Gustavo F. R. Gomes}
\end{large}
\end{center}

\vspace{0,5cm}
\begin{center}
\begin{normalsize}
\begin{large}
Orientator:
\end{large}
\end{normalsize}
\end{center}

\vspace{0.2cm}
\begin{center}
\begin{large}
\textbf{Professor Doutor Mário M. Freire}
\end{large}
\end{center}

\vspace{0,5cm}
\begin{center}
\begin{normalsize}
\begin{large}
Co-Orientator:
\end{large}
\end{normalsize}
\end{center}

\vspace{0.2cm}
\begin{center}
\begin{large}
\textbf{Professor Doutor Pedro R. M. Inácio}
\end{large}
\end{center}

\vspace{0.5cm}
\begin{center}
\begin{normalsize}
	June 8th, 2016
\end{normalsize}
\end{center}


\clearpage{\thispagestyle{empty}\cleardoublepage}

\frontmatter

\chapter*{Acknowledgements}
\label{chap:ack}

% A conclusão deste trabalho, bem como da grande maior parte da minha vida académica não seria possível sem a ajuda de ...
I would like to personally thank both my supervisors, Professor Mário M. Freire
and Professor Pedro R. M. Inácio, for the opportunity to develop this project
because it expanded my knowledge on various fields. I also want to acknowledge
their suuport and motivation throughout this endeavour.

Secondly I have to show my gratitude to all of my family, my parents, Fernando
Gomes and Eugénia Ramos, for educating me and supporting me, I would also like
to extend my gratitude to my grandparents, António Ramos and Maria Luísa
Veríssimo, and my uncles, Pedro Ramos and Sofia Ramos. Without my family I
would probably be a very different person.

I also want to recognize all the professors with whom I had the chance to learn
from, each and every one made me learn something useful that will influecne the
way I develop systems as well as the solutions I am able to design because of
what I have learned throughout my academic path.

Of course a honorable mention will go to Guido van Rossum, creator of the Python
programming language, not just because this project was developed with it but
also because he developed an entire sub-sulture which is based on open-source
and sharing knowledge without which I could achieve the system described in the
report.

Finally I want to give recognition to my friends who have supported and
encouraged me. Special mention to my best friend Filipa Costa for actually
understanding my dissapearance for along time.


\clearpage{\thispagestyle{empty}\cleardoublepage}


\tableofcontents

\clearpage{\thispagestyle{empty}\cleardoublepage}

\listoffigures

% Se não existirem tabelas, comentar as seguintes linhas
\clearpage{\thispagestyle{empty}\cleardoublepage}
\listoftables

% \listoflistings

\clearpage{\thispagestyle{empty}\cleardoublepage}
\chapter*{Acrónimos}
\label{chap:acro}

\begin{acronym}[TCP]
  \acro{TCP}{Transmission Control Protocol}
\end{acronym}

% \clearpage{\pagestyle{empty}\cleardoublepage}
% \chapter*{Glossário}
\makeglossaries

\newglossaryentry{.NET Framework}
{
  name={.NET Framework},
  description={É uma plataforma para desenvolvimento e funcionamento de aplicações desenvolvida pela Microsoft.}
}

\newglossaryentry{ASP.NET}
{
  name={ASP .Net},
  description={É uma plataforma da Microsoft para o desenvolvimento de aplicações Web e é o sucessor da tecnologia ASP.}
}

\newglossaryentry{CS}
{
  name={C\#},
  description={Lê-se \textit{C Sharp} e é uma linguagem de programação orientada a objectos, desenvolvida pela Microsoft, inicialmente para a plataforma .NET. O C\# é inspirado na junção entre as linguagens C++ e Java.}
}


\newglossaryentry{Java}
{
  name={JAVA},
  description={É uma linguagem de programação orientada a objectos, desenvolvida pela Sun Microsystems na década de 90. Hoje pertence à empresa Oracle.}
}


\newglossaryentry{OpenDMTP}
{
  name={OpenDMTP},
  description={\textit{Open Device Monitoring and Tracking Protocol} é um protocolo e uma \textit{framework} abertos que permite a comunicação bidireccional entre servidores e clientes através da internet.}
}


\newglossaryentry{OpenGTS}
{
  name={Open GTS},
  description={É o primeiro projecto \textit{Open Source} \textit{Web-Based} para controlo de frotas por GPS.}
}


\newglossaryentry{VS2010}
{
  name={Visual Studio 2010},
  description={\textit{Microsoft Visual Studio 2010} é um sistema de desenvolvimento desenvolvido pela Microsoft e é dedicado ao Framework .NET, que contem um conjunto de ferramentas de desenvolvimento projectadas para auxiliar os programadores a enfrentarem desafios complexos.}
}


\newglossaryentry{WebS}
{
	name={Web Service},
	description={Web services são aplicações modulares auto-descritas e auto-contidas, que permitem a integração de sistemas e a comunicação entre aplicações de diferentes tipos.}
}


\newglossaryentry{WebBased}
{
	name={Web Based},
	description={Aplicação desenvolvida para a Web.}
}

\newglossaryentry{Roaming}
{
	name={Roaming},
	description={Define a possibilidade de um utilizador de uma determinada rede obter rede/conecção fora da área geográfica onde foi registado.}
}


\newglossaryentry{Smartphone}
{
	name={Smartphone},
	description={Smartphone é um telefone móvel que contem muitas das principais tecnologias de comunicação e serviços que existem nos computadores pessoais, como acesso a e-mails, serviços de mensagens instantâneas, internet, GPS, entre outros.}
}

\newglossaryentry{TCPIP}
{
	name={TCP/IP},
	description={É um conjunto de protocolos de comunicação entre computadores ligados rede. O nome TCP/IP surge da união entre dois protocolos: o TCP (Transmission Control Protocol) e o protocolo IP (Internet Protocol).}
}

\newglossaryentry{Firewall}
{
	name={Firewall},
	description={É o nome criado para definir um dispositivo para uma rede de computadores que tem como objectivo criar uma política de segurança num determinado ponto de controlo da rede.}
}

\newglossaryentry{JavaScript}
{
	name={JavaScript},
	description={É uma linguagem de programação baseada na linguagem de programação ECMAScript. Actualmente é a linguagem de programação mais utilizada em \textit{``Client-Side''} nos \textit{browsers}.}
}

\newglossaryentry{Flash}
{
	name={Flash},
	description={Desenvolvido pela Macromedia, o Flash é um software utilizado para criação de animações interactivas que funcionam incorporadas em \textit{Browsers}, \textit{Desktop}, \textit{Smartphones}, \textit{Tablets}, e Televisores.}
}


\newglossaryentry{StoredProcedure}
{
	name={Stored Procedure },
	description={É o nome dado a um conjunto de comandos numa base de dados de forma a simplificar a sua utilização.}
}

\newglossaryentry{SQLS}
{
	name={SQL Server 2008},
	description={É um sistema de gestão de base de dados relacional criado pela Microsoft.}
}

\newglossaryentry{Firm}
{
	name={Firmware},
	description={É o conjunto de instruções operacionais programadas directamente no \textit{hardware} de um equipamento electrónico.}
}

\newglossaryentry{browser}
{
	name={Browser},
	description={É um programa de computador que possibilita aos utilizadores uma interacção com documentos virtuais da Internet, também conhecidos como páginas Web.}
}



\clearpage{\thispagestyle{empty}\cleardoublepage}

\mainmatter
\chapter{Introduction}
\label{chap:intro}

\section{Contextualization}
\label{sec:context} % CADA SECÇÃO DEVE TER UM LABEL
% CADA FIGURA DEVE TER UM LABEL
% CADA TABELA DEVE TER UM LABEL
We live in an information era, one that is driven by the use of more and more
devices connected to the Internet. These devices form the Internet of Things
(IoT). They could be a laptop, smartphone or even the new smart appliances in
our kitchen are connected to the Internet, which try to aid our everyday lives
by telling us that we ran out of milk or that there is strange movement around
house house even when you are at work.

% Os acrónimos devem ser definidos recorrendo ao pacote (\emph{package}) \texttt{acronym}, usando os comandos \texttt{\textbackslash acro}, \texttt{\textbackslash ac}, \texttt{\textbackslash acp}, etc. E.g., \emph{The subject of this report is network protocols, namely \ac{TCP}.  \ac{TCP} is studied for several aspects of performance.}


\section{Motivation}
\label{sec:mot}
Devices that are connected to the Internet are normally very limited when it
comes down to the computational resources at their disposal. This lack of
computational power leads to more relaxed security because those same security
measures that prevent intrusions or spying need some resources on their own.

As it can be expected big manufacturers of these devices won't care if their
product could be accessed through the Internet even by someone who's not their
client. This is very common unfortunately and there are more and more cases
reaching the public's eyes.

The challenge behind this project is to devise a way for any user be shielded
from the manufacturers negligent decisions and also notify the network user to
potential usage of unsecure protocols by their smart devices.

\section{Objectives}
% Goals would be better instead of objectives?
\label{sec:obj}
The project aims to protect every user's domestic network from potential
intrusions and avoid any privacy violations as well. The user could be
abstracted by the complexities of network configuration but at the same time
define exclusion or notifications for the usage of unsecure protocol.

\section{Distribution of Content}
\label{sec:organ}
% !POR EXEMPLO!
% De modo a refletir o trabalho que foi feito, este documento encontra-se estruturado da seguinte forma:
\begin{enumerate}
	% 1
	\item The first chapter -- \textbf{Introduction} -- talks about the context
		in wich the project will strive, it will also present the goals and the
		motives behind it's inception.
	% 2
	\item The second chapter -- \textbf{Related Works and Technology} --
		presents similar solutions to this particular project as well as the
		differentiationg factors between these solutions and this project at
		hand. It will also describe the technology used to conceive the system.
	% 3
	\item The third chapter -- \textbf{System Architecture} -- consists of a
		small requirement analysis, description of the architecture and also the
		tools used.
	% 4
	\item The forth chapter -- \textbf{System Description and Implementation} --
		is a description of all of its components as well as the implementation
		decisions involved in the creation of the subsystems.
	% 5
	\item The fifth chapter -- \textbf{Prototype and Testing} -- talks about the
		prototype as well as the setup in which the system was tested. It will
		also talk about how the system should be setup in real life, its
		installation and correct configuratin and usage.
	% 6
	\item The sixth chapter -- \textbf{Conclusions and Future Work} -- will
		consist in a description of the main conclusions extracted from the
		development of this system as well as what is going to be developed
		further in this system.
\end{enumerate}

%s
% \section{Algumas Dicas -- Retirar da Versão Final}
% ALGUMAS DICAS
% Os relatórios de projeto são individuais e preparados em \LaTeX, seguindo o formato disponível na página da unidade curricular. Deve ser prestada especial atenção aos seguintes pontos:
% \begin{enumerate}
%   \item O relatório deve ter um capítulo Introdução e Conclusões e Trabalho Futuro (ou só Conclusões);
%   \item A última secção do primeiro capítulo deve descrever suscintamente a organização do documento;
%   \item O relatório pode ser escrito em Língua Portuguesa ou Inglesa;
%   \item Todas as imagens ou tabelas devem ter legendas e ser referidas no texto (usando comando \texttt{\textbackslash ref\{\}}).
% \end{enumerate}

\clearpage{\thispagestyle{empty}\cleardoublepage}
\chapter{Related Works and Technology}
% Os titulos dados aos capítulos são meros exemplos. Cada relatório deve adequar-se ao projeto desenvolvido.
\label{chap:related-material}

\section{Introduction}
\label{chap2:sec:intro}
% Cada capítulo \underline{intermédio} deve começar com uma breve introdução onde é explicado com um pouco mais de detalhe qual é o tema deste capítulo, e como é que se encontra organizado (i.e., o que é que cada secção seguinte discute).
In this chapter there will be presented some related works, with some similar
functionalities as the project at hand. There will also be a section dedicated
to the description of the underlying techonolgy that is going to be used in the
development of this project.

\section{Similar Solutions}
\label{chap2:sec:similar-sol}
There are similar projects but not with the same set of functionalities as the
project at hand. The similar solutions are as follows:
\begin{enumerate}
	\item Nagios
	\item Domestic Routers
	\item Wireshark
	\item Snort
\end{enumerate}

Now we begin the description of the products listed above and also pointing out
the main differences between these solutions and the project described
throughout this report.

\subsection{Nagios}
\label{chap2:sec:similar-sol:sub:nagios}
One very similar project is the Nagios software, and it consists of a network
infrastructure monitoring system. Developed mostly for comercial and industrial
environments. There is a special version made for the Raspberry Pi, called
NagiosPi, that enables the user to check all of the connected devices and
monitor their traffic they generate. The NagiosPi system has the following
interfaces:
\begin{enumerate}
	\item Nagios;
	\item Nconf;
	\item NagVis;
	\item PHPMyAdmin;
	\item RaspControl;
\end{enumerate}

The NagiosPi system permits that the netwrok administrator logs all of the
information it might find pertinent as well as monitor all of the traffic
circulating inside the network. Using these interfaces the users can monitor
specific ports and hosts, it could visualize the packets and some statistics.

One of the biggest differences to the project being described in this report is
that it is more focused towards the domestic environments, because Nagios has a
complex interface using the web browser and not the mobile application offered
by this project. One more difference is that Nagios exclusively handles
monitoring, but the system being developed also has firewall capabilities, which
means it could block and also notification capabilities for flagged protocols
and packets.

\subsection{Domestic Routers}
\label{chap2:sec:similar-sol:sub:routers}
Another similar solution that could be referenced is the domestic router which
people have in their homes. This router actually comes with a firewall already
minimally configured preventing that any connection which isn't started by a
device inside the domestic network is unable to access or even probe the devices
which are inside the domestic network.

As said before this is a simple firewall but make a lot of difference because it
prevents malicious attackers from getting information about the network itself
and any of the devices connected to it, as well as prevents the exploitation of
these devices becasue they need that device to access a resource on the Internet
in order to become infected. This might sound useless but the reason wy people
get infected by a computer virus is because they access a website or resource on
the Internet which spreads malware, so it is clear that the user inside the
domestic network needs to access the infectious resource or download it in order
to get infected by it, in other words, it needs to bring the infection to
himself and not be always vulnerable.

The main difference between the firewall system of domestic routers and the
system described in this project is that the domestic router filters incoming
traffic, packets and connections coming from the Internet to the domestic
network. In the case of this system the filtering is done on the outgoins
packets, so they complement each other, preventing the user from getting malware
and also making it possible to prevent any communication between the virus and
any other resource on the Internet.
Another difference is that our system can be easily configurable using the
mobile application as well as notify the user of flagged behaviors or
communications, which the home router can't do.
% O trecho de código seguinte mostra a função \texttt{main()} e o seu funcionamento:
% \begin{lstlisting}[caption=Trecho de código usado no projeto.]
% #include <stdio.h>

% int main(){
%   int i = 0;
%   for(i = 0; i < 100; i++)
%     printf("%d\n",i);
% }
% \end{lstlisting}


% Se quiser definir a distribuição de Pareto, posso colocar a fórmula \emph{inline}, da seguinte forma $P(x)=\frac{x^{1/\Lambda}_{i}}{2}$, ou numa linha em separada, como se mostra a seguir:
% $$ y^2 = \sum_{x=0}^{20}( x^3 - 2x + 3).$$

% Outra maneira, mas numerada, é usar o ambiente \texttt{equation}, como se mostra na (\ref{eq:eq1}):
% \begin{equation}
%  y^2 = \sum_{x=0}^{20}( x^3 - 2x + 3).
%  \label{eq:eq1}
% \end{equation}

% \begin{align}
%  2+2+2+2+2+2+2+2+2+2+y^2 = & \sum_{x=0}^{20}( x^3 - 2x + 3);\\
%                          = & x^4 -2.
%  \label{eq:eq2}
% \end{align}

\section{Raspberry Pi}
\label{chap2:sec:rasp}
The Raspberry Pi is a small computer that comprises all of the vital component
that make it a computer, with a CPU, memory and input and output capabilities.
This tiny computer was choosen due to its affordable price, less then 45€. For
that price the resources available are not extremly performant but for the
development of this gateway it is enough due to the low-level nature of the data
being accessed.

This micro-computer will have installed the Operating System and also the
project itself which comprises a web service, a database, a firewall and a
notification scrip for emails.

%talk about the specs of the Raspberry Pi 2?

\section{Android}
\label{chap2:sec:android}
For the interaction between the user and system running on the Raspberry Pi an
Android application will be used. The Android OS is one of the major mobile
operating systems used and it is mostly due to the fact that it can run on
devices with low computational resources so it is the perfect platform for the
development of the companion app for this system.

This mobile operating system operates using applications written in Java which
is also one of the most popular programming languages in the world, the Android
OS is also equipped with a vast set of tools and libraries that can be leveraged
to create a web service consumer which makes it a perfect fit for the purposes
of this project.

% make a point that Android has inexpensive devices like smartphones and tables.

\section{Conclusions}
\label{chap2:sec:concs}
This chapter talked about some of the similar products and also the main
differentiating characteristics between these solutions and our project, like
the Nagios monitoring system, or the firewall of domestic routers.
In this chapter there was also a description of the main technologies being
leveraged in this project, the Raspberry Pi inexpensive computer and also the
Android mobile operating system.

\clearpage{\thispagestyle{empty}\cleardoublepage}
\chapter{Technologies and Tools Used}
% OU \chapter{Trabalhos Relacionados}
% OU \chapter{Engenharia de Software}
% OU \chapter{Tecnologias e Ferramentas Utilizadas}
\label{chap:tech}

% we could argue here the case for linux android and the raspberry pi

\section{Introduction}
\label{chap3:sec:intro}
This project will make use of the Netfilter subsystem of the Linux kernel, this
module will be used for the capturing of packets for further analysis. As for
the mobile app it will be implemented for the Android operating system and it
will make use of the Java programming language.

\section{Netfilter}
\label{chap3:sec:netfilter}
The Linux kernel has a framework that enables the interception and capture of
packets that are travelling through the network stack of any device that is
running Linux. The system is called Netfilter and permits the registration of
hooks at specific points in the network stack which allow an aplication to
manipulate and even to delete packets from the stack. This is a very powerful
framework that makes it possible to implement firewalls with less of a hassle.
The \emph{iptables} firewall system present in most Linux distribution uses this
system for its packet filtering capabilities, enabling the user of this
application to perform actions based on the hooks
% Talk about the fact that the hooks could be specified when filtering packets
% in the iptables POSTROUTING, PREROUTING, etc...

In the Linux kernel there is a framework capable of deploying hooks to the
network stack in order to manipuate the packets that are received. This
framework is used by \emph{iptables}, the firewall system that comes in all the
Linux distributions as default, to implement the rules of a firewall.

The hooks allow packets to be intercepted, while they are travelling through the
network stack of the device and these can be accessed in the following
instances:
% insert all of the hooks
\begin{enumerate}
	\item Pre-Routing
	\item Post-Routing
	\item Put them all here
\end{enumerate}
% Describe when one of those hooks occurs, what happens or where the packet is
% when the hooks is triggered

\section{GNU/Linux}
\label{chap3:sec:gnux}
GNU/Linux is the operating system environment that the system will run, it will
also make use of the Netfilter framework which is a component of the Linux
kernel. This operating system enables the use of less resources in comparison to
other Operating Systems available at the present time.
By using this OS the resources available for the application are sufficient even
for larger networks, another reason is the availability of low-level controls
ready for a developer to utilize, like the framework that enables the
implementation of the firewall capabilities of the project being developed.

Of course GNU/Linux is not the only operating system available with these
capabilities but it is the most adopted and also enables better performance then
most of its competitors.

\section{Raspberry Pi}
\label{chap3:sec:rasp}
The Raspberry Pi board is a small computer, in the sense that it has all of the
components of the von Neumann architecture, it has a Central Processing Unit
(CPU), memory and input and output slots for external devices.
For this project there was the use of a Raspberry Pi model 2, that boasts a
quad-core processor with a clock speed of 800MHz, one gigabyte of RAM (RANDOM
ACCESS MEMORY) and has a networking port and a microSD card slot. With these
features a Linux distribution is ran and it will act as a software firewall and
a default gateway giving it access to all the networks traffic, this way it can
fitler it according with the guidelines provided by the end-user of the system.

Choosing the Raspberry Pi as the device in which to run the system was simple.
Since this particular model and the new Raspberry Pi model 3 are available
between the price of 30 and 40€ it is an inexpensive machine that can function
as a gateway to outgoing traffic and anyone could buy one and install the system
in order to protect itself.

All of the system will be constrained by the lack of greater power but since the
main purpose of this device is to filter the outgoing traffic it will be
suficient even for bigger networks.

Tests should be done on both the Raspberry Pi 2 and 3 since they differ when it
comes to CPU performance. The main focus will be on the Raspberry Pi 2 but once
that is done a few tests should occur to check if it would be better performant
on the model 3.

\section{Packet Filtering System}
\label{chap3:sec:pfs}
The Packet Filtering System will be on the Raspberry Pi and will be responsible
for the capturing of all incoming packets to the system. It will make use of
the Netfilter susbsytem present in the Linux kernel, which willmake it possible
to implement hooks that will interact with the networking card of the Raspberry
Pi and it will enable the capture of every packet. This is one of the vital
parts of the system because of the efficiency gained from interacting with the
operating system's kernel directly.

By having access to every packet that arrives to the network card we can then
process them and rule according to the specified config file. For the filtering
per se we will use just plain C with arrays/lists  for the classification of
the protocols used.

\section{Android}
\label{chap3:sec:andrd}
Android  is an operating system developed for mobile platforms in particular for
smartphones and later also for tablets. The underlying technology was bough by
Google back in ......... and it has since evolved into a fully formed operating
system capabale of also working with limited resources present in mobile deices.
At it's heart it has a stripped down version of the Linux kernel, built to
provide good power to hardware performance.
% Talk about the wide use of the OS in mobile devices

By developing the companion application of this project in this operating system
two thing will be accomplished, there will be a bigger availability because of
the world wide adoption of this particular operating system and also it will
enable the end user to control the entire system with their mobile devices,
e.g., most people carry around their phone even when they are home so, by
providing a mobile application any user could access the system, and also
receive immediate notifications if something out of the ordinary is happening
inside their network.

This companion application will only work inside the network and it will
implement security mearsures in order to disable any connection from a possible
attacker. The only thing that could be risked is if the device that has the
companion application is compromised.

\section{Android Companion App}
\label{chap3:sec:aca}
The system will communicate securely with the a Android application, this
process will use communications through SSL/TLS and it will communicate with
the filtering system, which runs on the Raspberry Pi.

The application itself will act as a Command\&Control and it will abstract the
user from all the work being done within the Packet Filtering System.

This companion application will connect to the server-side and enable the
network administrator to categorize any protocol as unsafe and to choose which
action it should perform once a connection from a device is made to the
Internet. The user could also monitor the traffic and access some metrics from
the activity of the Local Area Network devices. These metrics could be simple
percentages of unsafe connections or even refering to a specific protocol. It
could also log how many hours and the amount of traffic to a specific host, it
could also list all of the hosts visited from the devices inside the domestic
network.

% A tabela~\ref{tab:exemplo} serve apenas o propósito da exemplificação de como se fazem tabelas em \LaTeX.
% \begin{table}
% \centering
% \begin{tabular}{|c|lr|}
% \hline
% \textbf{campo 1} & \textbf{campo 2} & \textbf{campo 3}\\
% \hline
% \hline
% 14 & 15 & 16 \\
% \hline
% 13 & 13 & 13 \\
% \hline
% \end{tabular}
% \caption{Esta é uma tabela de exemplo.}
% \label{tab:exemplo}
% \end{table}

\section{Conclusions}
\label{chap3:sec:concs}
The focus of this chaper was on what technologies and components involved in
this project.
The use of low-level systems will help efficiency wise. This will also help by
making the analysis process as soon as the packets arrive to the machine.
% Cada capítulo \underline{intermédio} deve referir o que demais importante se conclui desta parte do trabalho, de modo a fornecer a motivação para o capítulo ou passos seguintes.

\clearpage{\thispagestyle{empty}\cleardoublepage}
\chapter{System Description and Implementation}
% Os titulos dados aos capítulos são meros exemplos. Cada relatório deve adequar-se ao projeto desenvolvido.
\label{chap:sys-desc}

\section{Introduction}
\label{chap4:sec:intro}
% Cada capítulo \underline{intermédio} deve começar com uma breve introdução onde é explicado com um pouco mais de detalhe qual é o tema deste capítulo, e como é que se encontra organizado (i.e., o que é que cada secção seguinte discute).
This project contains a detailed description of each of the subsystem running on
the Raspberry Pi as well as the mobile application. As mentioned in the previous
chapter the system which will run on the Raspberry Pi is inspired by the concept
of Microservices meaning that each service provided by the gateway system is
independent from all of the other services, which means a detailed approach to
each service is necessary.

\section{Web Service Subsystem}
\label{chap4:sec:web-sys}

The Web Service is be responsible for all the communication between the user and
the Raspberry Pi components. The web service will be all implemented making use
of the \emph{Flask} framework allied with the \emph{flask-restful} extension
that enables the declaration of resources to be accessed through
\emph{\ac{HTTP}} requests.

In order to implement each resource we only need to create a class to represent
each \emhp{\ac{URI}}. These resources will handle the \emph{\ac{HTTP}} methods,
such as GET, POST, PUT and DELETE, according to their implementation providing
the following functionalities through the \emph{\ac{API}}:
\begin{enumerate}
	\item List the ruleset of the server;
	\item Add rules to the ruleset;
	\item Remove rules from the ruleset;
	\item Edit rules in the ruleset;
	\item Get the latest counter values for the entire ruleset;
	\item Get all of the values for a given rule.
\end{enumerate}

This Web Service implements a \emph{\ac{REST}} interface for the interaction
with external devices. The use of the \emph{\ac{REST}} \emph{\ac{API}} makes
interfacing with other devices more standard, opening the possibility to
implement other applications that access this \emph{\ac{API}}.
We now present the principles which need to be followed by the Web Service in
order to be considered a \emph{\ac{REST}} compliant web service:
\begin{enumerate}
	\item Client-Server architecture with a clear separation between the two.
	\item Stateless -- every request from the client must contain all the
		pertinent information, which includes authentication credentials, for
		example. The server should not store any session information about the
		client.
	\item Cache -- responses from the client might be chacheable or not by the
		server or intermidiate systems in order to maximing performance. These
		responses must be labeled by the server as cacheable or noncacheable.
	\item Uniform Interface -- the protocol used between the client and the
		server resources must be well defined and standardized, for exmaple the
		use of the HTTP protocol.
	\item Layered Interface -- other machines and systems might be added between
		the client and the server for reliability, scalability or performance
		reasons.
	\item Code-on-Demand -- optionally, clients can dowload and execute code
		from the server in their own context.
\end{enumerate}
(These principles were first described in Roy Fielding's Ph.D. dissertation.)

The Web Service will parse and respond with data in the \emph{\ac{JSON}} format
providing abstraction from the language being used for the client application.

The \emph{\ac{API}} for this RESTful web service can be accessed by appending
the following \emph{\ac{URI}} to the gateway's \emph{\ac{IP}} address:
\begin{enumerate}
	\item "rule/<int:id>" -- making it available for the client to check for a
		specific rule, as well as delete or edit it, using the GET, DELETE and
		PUT requests, respectively, from the \emph{\ac{HTTP}} protocol.
	\item "/rules" -- this \emph{\ac{URI}} will enable the client application to request for
		the full set of rules currently running on the server, by using the GET
		request. By making a request with the POST type the client application
		can append new rules to the rule set already in place, which will then
		be updated during the next cycle of the firewall subsystem.
	\item "/counter/<int:id>" -- URI that will only be accessed through GET
		request and provided with the id of the rule it will deliver the
		counters value, with number of packets and total number of bytes that
		were matched with that rule.
	\item "/counters" -- also only accessible through GET requests, which sends
		a list of the latest counters collected from the firewall service. These
		will be a list of rules identified by their ID, as well as the number of
		packets and bytes for that corresponding rule.
\end{enumerate}

% TALK ABOUT THE FIELDS OF THE flask-restful
In order to send all of the responses in a standard \emph{\ac{JSON}} format the
\emph{flask-restful} extension provides a formatter to convert the objects into
the \emph{\ac{JSON}} format intended. This is demonstrated in the following code
listing:

\begin{lstlisting}[caption=Example of the email JSON template.]
ip_fields = {
        'ip': fields.String,
        'ipv4': fields.Boolean,
        'src': fields.Boolean,
}

mac_fields = {
        'mac': fields.String,
}

rule_fields = {
        'id': fields.Integer,
        'port': fields.String,
        'protocol': fields.String,
        'action': fields.String,
        'ip': fields.Nested(ip_fields),
        'mac': fields.Nested(mac_fields),
}
\end{lstlisting}

This is then used by the \emph{marshal\_with} Python decorator in order to
convert the data object from the database into \emph{\ac{JSON}} format.

\begin{lstlisting}[caption=Example of the use of marshal\_with decorator.]
class RuleResource(Resource):
	# ...
	@marshal_with(rule_fields)
	def get(self, id):
		rule = session.query(Rule).filter(Rule.id == id).first()
		if not rule:
			abort(404, message="Can't find {0} in the DB.".format(id))
		return rule
	# ...
\end{lstlisting}

\section{Database Subsystem}
\label{chap4:sec:db-sys}
The database can be considered the center of the entire system since it will be
where all of the data stored to represent the firewall rules as well as the
packets which were flagged by the user. It will also contain pairs of username
and password representation for the user authentication even though this only
having one user registered might be enough for this type of system.
The database management system is SQLite 3 which as already a good API provided
in the Python standard libraries, nonetheless this project will make use of a
package called SQLAlchemy that can is freely available online which abstracts
the developer from the SQL queries.

The SQLAlchemy module enables the programmer to represt all of the entities as
well as the relationships between them using Python classes that inherit from
the Model superclass provided by SQLAlchemy. This way the developer only need to
declare the name of the attributes as well as the table name by declaring them
as attributes of the classes.

The models that need to be implemented in order for the whole system to work as
intended are as follows:
\begin{enumerate}
	\item User -- representing the user's credentials, username and password
		fields.
	\item Rule -- represents a rule for the firewall system to execute.
	\item Rule\_IP -- as the IP that the rule might have associated to it.
	\item Packet -- as a representation of a network packet, especially used for
		the notifications.
\end{enumerate}

These are all entities in the database and will conveil all of the data for the
application and the rest of the subsystems will all needs to routinely interact
with the database in order to get the rules for the firewall, to authenticate
users or even to get the information for the email and app notifications.

% insert and explain at least one of the models

After modelling the data we only need to create a session for the database, and
all of the data inside the database will be accessible with the call of methods
from the session instance. These methods could be insert, delete, update and
query.

% insert and explian at least one query and one insertion

% The database is the central component of our system and it is the most
% important system, because without it nothing else can run.

The database is organized with the following schema:

%insert an image of the database's schema here

% Describe each entity within the database.

\section{Email Notification Subsystem}
\label{chap4:sec:email-sys}
The Email Notification Subsystem will act as a complement for the Android
application, which means that it will only send out notifications to the user if
he isn't home. In order to detect the user, the Email Service queries the
database to check when was the last time the user was making requests to the Web
Service API, then calculates the time elapsed and if it's greater or equal to
the cycle period it will send an email with a small report about the network
traffic made inside the local network. It will then simply sleep until the next
cycle repeating the process. This cycle time will be defined by the user when he
initially configures the system in the Raspberry Pi.

The user will need to define the email address where he wants to receive the
notifications as well as the cycle periods, i.e., the time that needs to pass
between notifications sent to the users email, lastly the user will also be
asked to provide his first name in order for the system to address the user
properly.

The email being sent to the user is built using a template that is formatted
according to JSON standards which makes it highly customizable if the user
actually wants to modify some of the fields. The following code listing
demonstrates a typical template, the default template provided within the
installation files. During configuration, especially if the user uses the easy
configuration scrip some of these fields will suffer changes, both the
users\_email and the users\_name will be changed accordingly in this template
file after configuration.

The email service will also be in charge of creating the small reports that go
in the reason field, which is the body of the email per se. This is all done
automatically by the \emph{mixRuleCounters(. ,.)} which takes as parameters the
rules and counters received from the database. The method will then iterate
through all of the rules which have \emph{NOTIFY} or \emph{NOTIFY\_BLOCK} as
actions and link them with the counters for those rules.

Since the system needs to be clear for the user, all of the rules are explicitly
converted into a more user friendly format and then they will be united with the
values for the counters at a specific time.

% put a small template of the email as a code listing.

The emails being sent will all have the same design which is defined in the a
JSON formatted file named, \emph{email.json}. The Email Service will compose the
email for the user based on this template changing the various fields according
to the information that the user gave during initial configuration.

\begin{lstlisting}[caption=Example of the email JSON template.]
{
    "email": {
        "users\_name": "Gustavo Gomes",
        "sign": "Sent From Heimdall Secure Gateway.",
        "goodbye": "Thank you for your attention,",
        "subject": "Notification From Heimdall Secure Gateway",
        "reason": "Testing this email thingie",
        "email\_to": "gustavofrgomes@hotmail.com",
        "greeting": "Hello there Mr/Mrs"
    }
}
\end{lstlisting}

A problem found with this feature was that the email system globally has become
much more secure and it is impossible to send an email without an Email Provider
which will also imply that the user would need to purchase or get a domain name
in order to make this feature possible, the solution found for this problem was
to create an email account that will have all the credentials and be prepared to
be accessed by this subsystem, using the email address and the password of that
account and connecting via a TLS connection to the email provider, which in this
case is Google's Gmail. By doing this it is possible to send out emails to the
user's specified account from that account created for the system.

By making use of the \emph{email} module provided within Python's Standard
Library, it is possible to connect to an account on a SMTP server in order to
send any email, after the input of user's credentials such as the email address
and the password of that account.

% last thing!!!
% The Email Notification service will be keep notifying the user whenever he isn't
% at home, this is a measure of securing the local network, but it would also need
% a server to which both the gateway device and a possible mobile application
% would need to connect in order to make communication possible.

% Through email it will be easier for the device to simply send messages of what
% is happening with the devices inside their domestic network. If there are any
% flagged communitcations are performed then this subsystem will inform within
% regular, user defined, intervals of time, e.g. every two hours send a report if
% any activity, signalled by the user as note-worthy, has occured.

% As referred before this will be a complement for the Android notifications
% because this will only occur whenever the user is not home, this is all done by
% logging when the last request to the web service preventing that both the email
% and Android notifications to overlap.

% Using the \emph{email} module present in the Python standard library it was
% possible to connect to an SMTP server like Googles gmail and after providing
% access credentials, username and password, the system can now send emails for a
% specified email address.

\section{Firewall Interface Subsystem}
\label{chap4:sec:firewall-sys}
The Firewall Subsystem is the one component responsible for the management and
deployment of rules inside the firewall, iptables. The only concern for this
system is to deploy the rules from the database and the accounting for the
network traffic of those same rules. First we will talk about the
\emph{python-iptables}, or \emph{iptc} package, which enables the direct
interaction between this system and the firewall system, which was described in
a previous chapter, the \emph{iptables}. Then a description of the class that
manipulates and handles both the rules and their corresponding counters is
described as well as the processes involved with the data retrieval from the DB
and conversion.

\subsection{python-iptables}
\label{chap4:sec:firewall-sys:iptc}
The \emph{python-iptables} is a third party package that can be found in the
Python Package Index (PyPI), it provides a direct interface with the
\emph{iptables} firewall system. This is mostly due to the fact that it is a
Python wrapper for the \eph{libiptc} which, in turn is a C library that
provides a programming interface to the \emph{iptables}.

This module allows the access to any of the tables, chains and rules within the
\emph{iptables} system, as well as the addition of new rules to the system and
also the manipulation of existing ones. This means that it is possible to
dynnamically manipulate the rules inside the \emph{iptables} system in any way
the developer might require.
Some of the main functionalities that are provided by this module as well as
used in the context of this project are:
\begin{enumerate}
	\item Target declaration;
	\item Protocol definition;
	\item Port speficication;
	\item IP filtering, which includes IP ranges;
	\item MAC address filtering;
	% \item Pattern matching.
\end{enumerate}

The way these functionalities will be used is covered in the next chapter.

\subsection{Firewall Class}
\label{chap4:sec:firewall-sys:firewall-class}
There is a Firewall object that can be configured with the chains defined on
\emph{iptables}, this can be accomplished by setting defining the named optional
parameter when initializing a Firewall object, it also defaults to the 'FORWARD'
chain in the \emph{iptables} system which is the one that this subsystem needs
to interact with. It will always connect to the \emph{FILTER} table of
\emph{iptables} since this is the table that filters all of the traffic passing
through the machine, i.e. the Raspberry Pi.

Since the entire system is modular the Firewall class can be used independently
from the system itself. Of course it is still recommended that it should be
paired at least with the database since it has methods to query a database using
a \emph{SQLAlchemy} session.

The methods that are provided in this class are the following:
\begin{enumerate}
	\item addRule() -- method that handles the conversion of rules inside the
		database model object to \emph{python-iptables} rule objects.
	\item submitRule(.) -- handles the submission of the rule to the database,
		if the table from the database has the autocommit feature as False then
		it will simply add the rule but not commit it to the working firewall
		rule set.
	\item flushRules() -- responsible for flushing all of the rules from
		the chain.
	\item resetCounters() -- method that will reset the counters for all of the
		rules in the \emph{iptables}, as well as the data inside the database.
		This is used when the rules change because the counters need to start
		from zero and the database needs to be cleaned from the old rule set.
	\item ruleUpdate() -- main cycle for this subsystem, which gets the rules
		from the database, converts them from the models inside the database to
		\emph{iptables} compliant rules, resets the counters and adds the rules
		to the chain before ever commiting them to the firewall itself.
	\item createRule(.) -- creates an \emph{iptables} rule using a dictionary
		instead of an instance of the database models.
	\item counterUpdate() -- updates the counters inside the database, since the
		counters aren't dynamic per se, the table of the firewall needs to be
		refreshed in order to update the values of the counters. This will then
		publish the results in the database.
	\item counterDB(.) -- method that receives the results of the counters and
		connects them to each rule, afterwards sending it all to the database,
		by converting the data into a database model object.
	\item run(wait\_time) -- main method that will run both the ruleUpdate()
		method as well as the counterUpdate() to always keep the rules tracked
		and equal to the ones on the database. The \emph{wait\_time} param
		defaults to 30 and it is the periodicity of the cycle, in seconds.
\end{enumerate}

The Firewall Subsytem consists of the firewall itself, which will be
accomplished using the \emph{iptables} system, present in most Linux
distributions as the default firewall system, and there is also a firewall
handling system that is responsible for acquiring the rules, defined by the
user, from the database and deploying them on the firewall system.

\section{Android Companion App}
\label{chap4:sec:android}
In order to make it a complete system an Android mobile application was
conceived that mostly allows the user to interact with the Web Service's API.
The requests to the web service are made using the Volley as well as the
HTTPClient, since the Web Service only provides the API via HTTPS and with
authentication, the HTTPClient is required to handle the self-signed certificate
provided with the application as well as take care for the HTTPBasic
authentication headers.

As already said when describing the Web Service all of the data transmitted over
the network needs to be in the JSON format in order to provide a standard way of
communicating with all the client devices.

The aplication will irst ask the user to provide the authentication credentials
as well as the gateway's IP address.

% insert the authentication print

There are some classes defined to help structure all of the data and to allow
for some organization and manipulation of the data that can be received by the
Web Service. These classes are the Rule, Ip and Mac classes.

The activities that are provided by the application are presented in the
following list:
\begin{enumerate}
	\item AddRule -- Activity to append a rule at the end of the rule set on the
		server, uses the POST request;
	\item EditRule -- Used to edit a rule that is already implemented, making
		use of the PUT request method.
	\item DeleteRule -- Small activity that handles the DELETE request for a
		specific rule;
	\item ListRules -- Activity to list and access all of the rules that are in
		the system.
	\item Login -- First activity to be launched providing the user with a way
		to input both the credentials for authentication in the Web Service as
		well as the IP and Port number of the gateway.
	\item MainMenu -- This activity is just a simple menu to enable the user to
		interact with all of the functionalities proviced by the application.
\end{enumerate}

Almost all of the activities make use of the Volley package, described in
Chapter 2, for the communication between the Web Service and the application.
The application composes Hyper Text Tranfer Protocol (HTTP) requests to be sent
to the Web Server as well as JSONObject or JSONArray objects to send and parse
information received from the Web Service, running on the gateway device.

\section{Conclusions}
\label{chap4:sec:concs}
This chapter described all of the subsystems that comprise the project. We
talked about the Web Service with the Flask server and also the Android app that
consumes it. We also described the Database system which is the central
component and all systems depend on it. We referred the Email Notification and
the details that need to be in place to make email notifications possible. At
last we described how the Firewall subsystem and how we accomplished the sense
of dynamic firewall configuring.

\clearpage{\thispagestyle{empty}\cleardoublepage}
\chapter{Conclusões e Trabalho Futuro}
\label{chap:conc-trab-futuro}

\section{Conclusões Principais}
\label{sec:conc-princ}

Esta secção contém a resposta à questão: \\
\emph{Quais foram as conclusões princípais a que o(a) aluno(a) chegou no fim deste trabalho?}

\section{Trabalho Futuro}
\label{sec:trab-futuro}

Esta secção responde a questões como:\\
\emph{O que é que ficou por fazer, e porque?}\\
\emph{O que é que seria interessante fazer, mas não foi feito por não ser exatamente o objetivo deste trabalho?}\\
\emph{Em que outros casos ou situações ou cenários -- que não foram estudados no contexto deste projeto por não ser seu objetivo -- é que o trabalho aqui descrito pode ter aplicações interessantes e porque?}
\clearpage{\thispagestyle{empty}\cleardoublepage}

% SE EXISTIREM APENDICES, DESCOMENTAR O QUE ESTÁ EM BAIXO
% \appendix
% \include{apendice1}
% \clearpage{\pagestyle{empty}\cleardoublepage}
% \include{continuacao}
% \clearpage{\pagestyle{empty}\cleardoublepage}
% \include{apendice2}
% \clearpage{\pagestyle{empty}\cleardoublepage}
% \include{apendice3}
% \clearpage{\pagestyle{empty}\cleardoublepage}

\backmatter

\bibliographystyle{plain}
\thebibliography{bibliografia}

\end{document}